\documentclass{article}
\usepackage{amsmath}
\usepackage{amssymb}
\usepackage{amsthm}
\usepackage{MnSymbol}
\usepackage{xcolor}
\usepackage{parskip}
\usepackage{tikz}
\usepackage{bbm}
\usepackage{hyperref}
\usetikzlibrary{trees}
\usetikzlibrary{graphs}


\theoremstyle{definition}
\newtheorem{definition}{Definition}
\newtheorem{theorem}{Theorem}


\newcommand{\abs}[1]{\left| #1 \right|}


\title{lecture}
\date{\today}
\begin{document}
\maketitle

at dimension n, we can embed any ordinal relationship

\(d\leq n\) embedding is always possible

use matrix with epsilons, row vectors is embedding for that entry. computing shows this embedding respects ordinality

\begin{equation}
    \frac{n}{2}\leq d \leq n
\end{equation}


lower bound is proven by saying there are too many comparisons (asymptotically too many), and there arent enough comparisons in \(\mathbb R^{n/2}\)


the polynomials have degree 2, since the eucl norm squared gives just \(x_1^2+\dots x_d^2\). there are \(n\cdot d\) variables, and there are \({{n \choose 2}\choose 2 }\approx n^4\) polynomials.


\section{Terminal (VIP) ordinal embeddings}

\begin{equation}
    T=\{t_1\dots t_k\}, V\setminus T = \{v_1\dots v_{n-k}\}
\end{equation}

you only care about comparing distances of only the terminal nodes, i.e. of the form
\begin{equation}
    (t,.)<(t^\prime,.)
\end{equation}

comparisons between \(v\)s is irrelevant


\textbf{this can be done in just \(k\) dimensions}


assign all \(t_i\) to \(\vec e_i \cdot k^3n^2\)

the vertex \(t\in T\), \(v\in V\), rank \(r(t,v)\in \{1,2,\dots k\cdot n\}\), specifying the order in which distance \((t,v)\) appears w.r.t. other vertecies. 

to compare \((t,v)\) and \(t^\prime, v^\prime\), 
\begin{equation}
    \abs{\varphi(t)-\varphi(v)}= \sum_{i\neq t}^{}r(i,v)^2 + \left( (r(t,v)+M) \right)^2 = \sum_{i}^{}r^2(i,v)+2r(t,v)M+M^2
\end{equation}

but the first term is  \(\leq k\cdot (km)^2\leq M\)



\section{diameter}

diameter of a set of points is max 1 norm between two points. 


n is number of points and \(d\) is the dimension
naive is \(O(n^2d)\), better is \(O(nd2^d)\)

generate \(2^d\) sign vectors \(s\) of form \(\prod_d \pm 1 \)

make \(f\) such that \(f(p_i) = s_i\cdot p\)

\(f\) over all p has \(n 2^d\) entries. 

\begin{equation}
    \abs{p-q}_1 = \abs{f(p)-f(q)}_\infty
\end{equation}












\end{document}